\documentclass[11pt,a4paper]{article}
\usepackage[utf8]{inputenc}
\usepackage[a4paper,top=2.3cm,bottom=2.6cm,left=2.5cm,right=2.5cm,headsep=0.6cm,footskip=1.1cm]{geometry}
\usepackage{amsmath,amssymb}
\usepackage{enumitem}
\usepackage{fancyhdr}
\usepackage{titlesec}
\usepackage{xcolor}
\titlespacing*{\section}{0pt}{1.3em}{0.2em}
\usepackage{hyperref}
\hypersetup{colorlinks=true,linkcolor=black,urlcolor=black,
  pdftitle={FZT 203E Problem Set 2: Solutions},pdfauthor={Mehmet Özkan}}

\newcommand{\Z}{\mathbb{Z}}
\newcommand{\R}{\mathbb{R}}
\newcommand{\C}{\mathbb{C}}
\newcommand{\gen}[1]{\langle #1\rangle}
\DeclareMathOperator{\ord}{ord}
\DeclareMathOperator{\sgn}{sgn}
\DeclareMathOperator{\Cl}{Cl}
\DeclareMathOperator{\lcm}{lcm}

\pagestyle{fancy}
\fancyhf{}
\renewcommand{\headrulewidth}{0pt}
\fancyfoot[L]{\small FZT 203E $\cdot$ Problem Set 2 $\cdot$ Solutions}
\fancyfoot[R]{\small\thepage}

\setlength{\parindent}{0pt}
\setlength{\parskip}{0.45em}
\setlength{\emergencystretch}{1.5em}

% The web page teaching/group-theory/problem-sets/ps2-solutions.html is generated
% from the part between the SOLUTIONS markers; keep to these commands there:
% \prob, parts, \pq, \ans, itemize/\item, \emph, $...$, \[...\], -- and blank lines.
\newcommand{\prob}[1]{\par\vspace{1.1em}\noindent{\bfseries\large #1}\par\vspace{0.1em}}
\newenvironment{parts}{\begin{enumerate}[label=\textbf{(\alph*)},leftmargin=2.2em,itemsep=0.6em,topsep=0.2em]}{\end{enumerate}}
\newcommand{\pq}[1]{\item {\color{black!60}#1}\par}
\newcommand{\ans}[1]{{\boldmath\bfseries #1}}
\setlist[itemize]{label=$\bullet$,itemsep=0.15em,topsep=0.1em,leftmargin=1.4em}

\begin{document}

{\small FZT 203E $\cdot$ Theory of Groups in Physics $\cdot$ Fall 2026 \hfill Istanbul Technical University\\
Instructor: Prof.\ Mehmet Özkan}

\vspace{0.9em}
{\LARGE\bfseries Problem Set 2: Solutions}

\vspace{0.5em}
Each part repeats the question in grey, then answers it; the results asked for are in bold. Conventions as in the problem set: products are read right to left, so in $gh$ the element $h$ acts first, and a permutation sends position $i$ to position $j$.

%BEGIN-SOLUTIONS
\prob{1. Direct products, lattices, parity and time reversal}
\begin{parts}
\pq{The order of $(g, k)$ in $G \times K$.}
In $G \times K$ the product is taken slot by slot, so $(g, k)^m = (g^m, k^m)$. This is the identity $(e, e)$ exactly when $g^m = e$ and $k^m = e$, that is, when $m$ is a multiple of both $\ord(g)$ and $\ord(k)$. The smallest such $m > 0$ is their least common multiple: \ans{$\ord(g, k) = \lcm\bigl(\ord(g), \ord(k)\bigr)$}.

\pq{When is $\Z_m \times \Z_n$ cyclic?}
If $\gcd(m, n) = 1$, then $\lcm(m, n) = mn$, so by (a) the element $(1, 1)$ has order $mn$, the order of the whole group. It generates $\Z_m \times \Z_n$, and $k \mapsto k\,(1, 1)$ is an isomorphism: \ans{$\Z_m \times \Z_n \cong \Z_{mn}$}.

If $\gcd(m, n) > 1$, every element $(a, b)$ satisfies $\lcm(m, n)\,(a, b) = (0, 0)$, so its order divides $\lcm(m, n) = mn/\gcd(m, n) < mn$. No element has order $mn$, so \ans{$\Z_m \times \Z_n$ is not cyclic}.

\pq{How many sites does the particle visit?}
After $t$ steps the particle sits at $t\,(a, b)$, so it visits exactly the cyclic subgroup generated by $(a, b)$: the number of sites is the order of $(a, b)$. In $\Z_6$ the element $a$ has order $6/\gcd(a, 6)$, in $\Z_9$ the element $b$ has order $9/\gcd(b, 9)$, and (a) combines them.
\begin{itemize}
  \item $(1, 1)$: $\lcm(6, 9) = $ \ans{18 sites}.
  \item $(2, 3)$: $2$ has order 3 in $\Z_6$ and $3$ has order 3 in $\Z_9$: \ans{3 sites}.
  \item $(3, 1)$: orders 2 and 9: \ans{18 sites}.
\end{itemize}
Every order divides $\lcm(6, 9) = 18$, so \ans{the largest number is 18}, and \ans{the particle never visits all 54 sites}: $\gcd(6, 9) = 3$, so $\Z_6 \times \Z_9$ is not cyclic. On the $4 \times 9$ lattice $\gcd(4, 9) = 1$, so $\Z_4 \times \Z_9 \cong \Z_{36}$, and \ans{$(1, 1)$ visits all 36 sites}.

\pq{Parity and time reversal form $\Z_2 \times \Z_2$.}
$(P^2 x)(t) = -(-x(t)) = x(t)$ and $(T^2 x)(t) = x(-(-t)) = x(t)$, so $P^2 = T^2 = e$. Also $(PTx)(t) = -(Tx)(t) = -x(-t)$ and $(TPx)(t) = (Px)(-t) = -x(-t)$, so $PT = TP$. The four maps $e, P, T, PT$ are closed under composition, and every one of them squares to $e$, since $(PT)^2 = P^2 T^2 = e$. A group of order 4 in which no element has order 4 is not $\Z_4$: \ans{$\{e, P, T, PT\} \cong \Z_2 \times \Z_2$}.

\pq{Which of $P$, $T$, $PT$ are symmetries?}
If two maps send solutions to solutions, so does their product; $e$ does; and each element is its own inverse. So the symmetries form a subgroup. Now let $x(t)$ solve $m\ddot x = F(x, \dot x)$.
\begin{itemize}
  \item $y = Px = -x$ has $m\ddot y = -F(x, \dot x)$, which must equal $F(y, \dot y) = F(-x, -\dot x)$: $P$ is a symmetry exactly when $F(-x, -v) = -F(x, v)$.
  \item $y(t) = x(-t)$ has $\dot y(t) = -\dot x(-t)$ and $\ddot y(t) = \ddot x(-t)$, so $m\ddot y(t) = F\bigl(x(-t), \dot x(-t)\bigr)$, which must equal $F\bigl(x(-t), -\dot x(-t)\bigr)$: $T$ is a symmetry exactly when $F(x, -v) = F(x, v)$.
  \item $y(t) = -x(-t)$ has $\dot y(t) = \dot x(-t)$ and $\ddot y(t) = -\ddot x(-t)$: $PT$ is a symmetry exactly when $F(-x, v) = -F(x, v)$.
\end{itemize}
\begin{itemize}
  \item Spring, $F = -kx$: all three conditions hold, \ans{$\{e, P, T, PT\}$}.
  \item Spring with friction, $F = -kx - \gamma v$: odd under $(x, v) \to (-x, -v)$, but not even in $v$, \ans{$\{e, P\}$}.
  \item Falling ball, $F = -mg$: even in $v$, but $F(-x, -v) = -mg \ne mg$, \ans{$\{e, T\}$}.
  \item Ball with drag, $F = -mg - \gamma v$: none of the three conditions holds, \ans{$\{e\}$}.
\end{itemize}
\ans{A force that singles out a direction in space, such as uniform gravity, destroys $P$; a force odd in the velocity, such as friction or drag, which dissipates energy, destroys $T$.}
\end{parts}

\prob{2. The symmetries of benzene}
\begin{parts}
\pq{Conjugating reflections; the conjugacy classes of $D_6$.}
With $s r^j = r^{-j} s$: $r^{-j}(r^k s)\, r^j = r^{k - j} s r^j = r^{k - j} r^{-j} s = r^{k - 2j} s$, and $s^{-1}(r^k s)\, s = s r^k = r^{-k} s$.

Rotations: $r^{-j} r^k r^j = r^k$ and $s^{-1} r^k s = r^{-k}$, so the class of $r^k$ is $\{r^k, r^{-k}\}$. Reflections: conjugation changes the exponent $k$ of $r^k s$ by an even number or reverses its sign, which keeps $k$ even or odd, and every exponent of the same parity is reached. The classes are
\[
  \begin{gathered}
  \{e\}, \quad \{r^3\}, \quad \{r, r^5\}, \quad \{r^2, r^4\}, \\
  \{s, r^2 s, r^4 s\}, \quad \{rs, r^3 s, r^5 s\},
  \end{gathered}
\]
with sizes \ans{$1 + 1 + 2 + 2 + 3 + 3 = 12$}.

\pq{The six reflections.}
\begin{itemize}
  \item $s = (26)(35)$: mirror through atoms 1 and 4.
  \item $rs = (12)(36)(45)$: through the midpoints of the bonds 1--2 and 4--5.
  \item $r^2 s = (13)(46)$: through atoms 2 and 5.
  \item $r^3 s = (14)(23)(56)$: through the midpoints of the bonds 2--3 and 5--6.
  \item $r^4 s = (15)(24)$: through atoms 3 and 6.
  \item $r^5 s = (16)(25)(34)$: through the midpoints of the bonds 1--6 and 3--4.
\end{itemize}
\ans{The mirrors of $s, r^2 s, r^4 s$ pass through two opposite atoms (each fixes two atoms); those of $rs, r^3 s, r^5 s$ pass through the midpoints of two opposite bonds (they fix no atom).} A symmetry of the hexagon carries a mirror through atoms to a mirror through atoms, so the two kinds cannot be conjugate.

\pq{$Z$, $N$, $H$ are normal, $\{e, s\}$ is not; $r^3$ is central.}
Each of $Z$, $N$, $H$ is closed under products (in $H$: $r^{2a} \cdot r^{2b} s = r^{2(a+b)} s$ and $r^{2a} s \cdot r^{2b} s = r^{2(a - b)}$), hence a subgroup. A subgroup is normal exactly when it is a union of conjugacy classes, and by (a)
\[
  \begin{gathered}
  Z = \{e\} \cup \{r^3\}, \\
  N = \{e\} \cup \{r^2, r^4\}, \\
  H = \{e\} \cup \{r^2, r^4\} \cup \{s, r^2 s, r^4 s\}.
  \end{gathered}
\]
\ans{$\{e, s\}$ is not normal}: it contains $s$ but not $r^{-1} s r = r^4 s$. Rotations commute with $r^3$, and $s r^3 = r^{-3} s = r^3 s$, so \ans{$r^3$ commutes with every element}. It is the rotation by $180^\circ$, \ans{$(x, y) \mapsto (-x, -y)$, the inversion through the centre of the ring}.

\pq{$H \cong D_3$, $D_6 \cong D_3 \times \Z_2$ and $D_6/Z \cong D_3$.}
$r^2 = (135)(246)$ and $s = (26)(35)$ map $\{1, 3, 5\}$ to itself, so the group $H$ they generate does too. The other six elements carry atom 1 to an even atom: $r$, $r^3$, $r^5$ send 1 to 2, 4, 6, and so do $rs$, $r^3 s$, $r^5 s$ by (b). So \ans{$H$ is exactly the set of symmetries that map the triangle $1, 3, 5$ to itself}. On this triangle $r^2$ acts as the rotation by $120^\circ$ and $s$ as the reflection through atom 1, and only $e$ fixes all three atoms, so \ans{$H \cong D_3$}.

$H \cap Z = \{e\}$ because $r^3 \notin H$. If $hz = h'z'$, then $h'^{-1} h = z' z^{-1} \in H \cap Z$, so $h = h'$ and $z = z'$. The $6 \cdot 2 = 12$ products $hz$ are therefore distinct and make up all of $D_6$. Since $z$ commutes with everything, $(hz)(h'z') = (hh')(zz')$, so $(h, z) \mapsto hz$ is a bijective homomorphism: \ans{$H \times Z \cong D_6$, that is, $D_6 \cong D_3 \times \Z_2$}. The map $hz \mapsto h$ is a homomorphism of $D_6$ onto $H$ with kernel $Z$, so by the first isomorphism theorem \ans{$D_6/Z \cong D_3$}.

\pq{The quotient $D_6/N$.}
The cosets are $N = \{e, r^2, r^4\}$, $rN = \{r, r^3, r^5\}$, $sN = \{s, r^2 s, r^4 s\}$ and $rsN = \{rs, r^3 s, r^5 s\}$. Their squares are $(rN)^2 = r^2 N = N$, $(sN)^2 = N$ and $(rsN)^2 = (rs)^2 N = N$. Every element of the quotient squares to the identity, so \ans{$D_6/N \cong \Z_2 \times \Z_2$}, not $\Z_4$. The two questions: \ans{is it a rotation or a reflection, and does it map the triangle $1, 3, 5$ to itself or exchange it with the triangle $2, 4, 6$?} $N$: rotation, keeps; $rN$: rotation, exchanges; $sN$: reflection, keeps; $rsN$: reflection, exchanges.
\end{parts}

\prob{3. Angles, spokes and the wagon-wheel effect}
\begin{parts}
\pq{$\R/2\pi\Z \cong U(1)$.}
$E(\theta + \varphi) = e^{i\theta} e^{i\varphi} = E(\theta) E(\varphi)$, and every $z$ with $|z| = 1$ is $e^{i\theta}$ for some $\theta$, so $E$ is a homomorphism onto $U(1)$. $e^{i\theta} = 1$ exactly when $\cos\theta = 1$ and $\sin\theta = 0$, that is, $\theta \in 2\pi\Z$: \ans{$\ker E = 2\pi\Z$}, and the first isomorphism theorem gives \ans{$\R/2\pi\Z \cong U(1)$}. The intervals $(-\pi + 2\pi k, \pi + 2\pi k]$, $k \in \Z$, cover $\R$ without overlapping, so $\theta$ lies in exactly one of them, and exactly one element $\theta - 2\pi k$ of the coset $\theta + 2\pi\Z$ lies in $(-\pi, \pi]$.

\pq{The wheel with $n$ spokes.}
$\theta \mapsto e^{in\theta}$ is multiplication by $n$ followed by $E$, a homomorphism onto $U(1)$. Its kernel is $\{\theta : n\theta \in 2\pi\Z\} = (2\pi/n)\Z$. The intervals of length $2\pi/n$ give the representative in $(-\pi/n, \pi/n]$ as in (a).

\pq{The stagecoach on film.}
$\omega = v/(0.5\,\text{m}) = 2v$ in rad/s for $v$ in m/s, so between frames the wheel turns by $\omega/24 = v/12$ radians. With 12 spokes the film records $v/12$ modulo $2\pi/12 = \pi/6$, and the audience sees the representative in $(-\pi/12, \pi/12]$.
\begin{itemize}
  \item The film shows the true rotation while $v/12 \le \pi/12$: \ans{up to $v = \pi$\,m/s $\approx 3.1$\,m/s}, about 11\,km/h.
  \item The wheel seems to stand still when $v/12$ is a multiple of $\pi/6$: \ans{$v = 2\pi k$\,m/s, $k = 1, 2, \dots$}, that is $6.3$, $12.6$, $18.8$\,m/s, \dots
  \item $v = 5$\,m/s: the true turn is $5/12 \approx 0.417$\,rad per frame; subtracting $\pi/6 \approx 0.524$ gives $-0.107$\,rad, which lies in $(-\pi/12, \pi/12]$. The apparent angular velocity is $24\,(5/12 - \pi/6) = 10 - 4\pi \approx -2.6$\,rad/s: \ans{the wheel seems to turn slowly backwards, at about 2.6\,rad/s} (0.4 turns per second), although it turns forwards at 10\,rad/s.
  \item $v = 7$\,m/s: $7/12 - \pi/6 \approx 0.060$\,rad per frame, so the apparent angular velocity is $14 - 4\pi \approx 1.4$\,rad/s: \ans{slowly forwards, at about 1.4\,rad/s}, although it turns at 14\,rad/s.
\end{itemize}

\pq{Lattice waves and the Brillouin zone.}
$e^{ikja} = e^{ik'ja}$ for every $j$ exactly when $e^{i(k - k')ja} = 1$ for every $j$. For $j = 1$ this says $(k - k')a \in 2\pi\Z$, and then it holds for every $j$. So the phases agree at every atom \ans{exactly when $k - k' \in (2\pi/a)\Z$}. For $k = 7\pi/(2a)$, subtract $2 \cdot 2\pi/a = 8\pi/(2a)$: \ans{$k' = -\pi/(2a)$}.
\end{parts}

\prob{4. Shuffles, recurrence and the 15 puzzle}
\begin{parts}
\pq{The two shuffles.}
Write $i \to j$ when the card at position $i$ moves to position $j$. In the out-shuffle the cards from positions $1, 5, 2, 6, 3, 7, 4, 8$ land in positions $1, \dots, 8$, so $1 \to 1$, $5 \to 2$, $2 \to 3$, $6 \to 4$, $3 \to 5$, $7 \to 6$, $4 \to 7$, $8 \to 8$:
\[
  \begin{gathered}
  \text{out-shuffle} = (235)(476), \\
  \text{in-shuffle} = (124875)(36).
  \end{gathered}
\]
For the in-shuffle, $5 \to 1$, $1 \to 2$, $6 \to 3$, $2 \to 4$, $7 \to 5$, $3 \to 6$, $8 \to 7$, $4 \to 8$. The order of a permutation is the least common multiple of its cycle lengths: \ans{3 out-shuffles} and \ans{6 in-shuffles} restore the deck. A cycle of length $\ell$ has sign $(-1)^{\ell - 1}$, so the out-shuffle has sign $(+1)(+1) = $ \ans{$+1$} and the in-shuffle $(-1)(-1) = $ \ans{$+1$}: both are even.

\pq{Recurrence.}
After $m$ steps the configuration is $p^m$ applied to the start. $S_N$ is finite, so $p$ has a finite order $k$, $p^k = e$, and the start recurs after $k$ steps. Since $p^m = e$ exactly when $\ord(p)$ divides $m$, \ans{the first recurrence is after $\ord(p)$ steps}. A particle on a cycle of length $\ell$ moves round that cycle and \ans{first returns after $\ell$ steps}. The order is the least common multiple of the cycle lengths, which can exceed $N$. Among the cycle types of $S_7$ the largest least common multiple comes from $3 + 4$: \ans{the largest order in $S_7$ is 12}, for example $(123)(4567)$. In $S_{10}$, \ans{$(12)(345)(6\,7\,8\,9\,10)$} has order $\lcm(2, 3, 5) = 30$.

\pq{The 15 puzzle: moves and parity.}
A move exchanges a tile with the empty square, the sixteenth tile, and moves nothing else: \ans{a transposition}. Colour the squares like a chessboard. Each move takes the empty square to a neighbouring square, of the other colour, so it is back on its starting colour, in particular on square 16, only after an even number of moves. The arrangement is then a product of an even number of transpositions that fixes 16: \ans{an even permutation of the tiles $1, \dots, 15$}.

\pq{Two arrangements.}
Exchanging only 14 and 15 is a single transposition, an odd permutation: \ans{not reachable}. Reading $15, 14, \dots, 1$ is the permutation $i \mapsto 16 - i$ of the tiles,
\[
  \begin{gathered}
  (1\,15)(2\,14)(3\,13)(4\,12) \\
  (5\,11)(6\,10)(7\,9),
  \end{gathered}
\]
with 8 fixed: seven transpositions, odd, so \ans{not reachable either}.
\end{parts}

\prob{5. Rotations of a cube}
\begin{parts}
\pq{The 24 rotations.}
\begin{itemize}
  \item Through the centres of opposite faces: 3 axes, turns by $90^\circ$, $180^\circ$, $270^\circ$: 9 rotations.
  \item Through opposite vertices (the body diagonals): 4 axes, turns by $120^\circ$, $240^\circ$: 8 rotations.
  \item Through the midpoints of opposite edges: 6 axes, turns by $180^\circ$: 6 rotations.
\end{itemize}
With $e$: \ans{$1 + 9 + 8 + 6 = 24$}. There can be no more: a rotation is fixed by the face it moves to the top (6 choices) and the turn of that face (4 choices), at most $6 \cdot 4 = 24$.

\pq{The permutations of the diagonals.}
If $h$ carries $d_i$ to $d_j$ and $g$ carries $d_j$ to $d_k$, then $gh$ carries $d_i$ to $d_k$: $\varphi(gh) = \varphi(g)\varphi(h)$.
\begin{itemize}
  \item $R$: $(1, 1, 1) \mapsto (-1, 1, 1) \mapsto (-1, -1, 1) \mapsto (1, -1, 1) \mapsto (1, 1, 1)$, so \ans{$\varphi(R) = (1234)$}.
  \item $C$ fixes $(1, 1, 1)$ and sends $(-1, 1, 1) \mapsto (1, -1, 1) \in d_4$, $(1, -1, 1) \mapsto (1, 1, -1) \in d_3$, $(-1, -1, 1) \mapsto (1, -1, -1) \in d_2$: \ans{$\varphi(C) = (243)$}.
  \item $B$ sends $(1, 1, 1) \mapsto (1, 1, -1) \in d_3$ and $(-1, -1, 1) \mapsto (-1, -1, -1) \in d_1$, and maps $d_2$ and $d_4$ to themselves: \ans{$\varphi(B) = (13)$}.
\end{itemize}

\pq{$\varphi$ is an isomorphism.}
If $g$ maps every diagonal to itself, it sends each vertex $v$ to $v$ or to $-v$. Neighbouring vertices are at distance 2, for example $v = (1, 1, 1)$ and $w = (-1, 1, 1)$. If $gv = v$ and $gw = -w$, the distance would become $|v + w| = |(0, 2, 2)| = 2\sqrt2 \ne 2$, but rotations preserve distances. So neighbouring vertices keep the same sign, and since the edges connect all eight vertices, $g = I$ or $g = -I$. The inversion $-I$ has determinant $-1$ and is not a rotation, so \ans{$\ker\varphi = \{e\}$}. $\varphi$ is injective and $|G| = |S_4| = 24$, so \ans{$\varphi$ is an isomorphism, $G \cong S_4$}.

\pq{Kinds of rotations and conjugacy classes.}
Orders are preserved. The 6 quarter-turns (order 4) go to the 6 four-cycles, and the 8 turns by $\pm 120^\circ$ (order 3) to the 8 three-cycles. The 9 half-turns go to the 9 elements of order 2: $\varphi(R^2) = (1234)^2 = (13)(24)$ is a double transposition, and so are the images of the other two face half-turns, which are conjugate to $R^2$ by $C$; this uses up the 3 double transpositions, so the 6 edge half-turns, such as $B$ with $\varphi(B) = (13)$, go to the 6 transpositions. The counts agree: $1 + 6 + 3 + 8 + 6 = 24$. In $S_4$ the classes are the cycle types, and an isomorphism maps classes to classes, so \ans{two rotations are conjugate exactly when they are of the same kind}.

For $g\colon (x, y, z) \mapsto (x, -y, -z)$, the half-turn about the $x$-axis, $g^{-1} = g$ and $g^{-1} R g\,(x, y, z) = g\,R(x, -y, -z) = g(y, x, -z) = (y, -x, z) = R^{-1}(x, y, z)$. So \ans{$g^{-1} R g = R^{-1}$}. Physically, $g$ turns the $z$-axis upside down, and a counterclockwise quarter-turn seen from the other end of its axis is clockwise: \ans{the turns by $+90^\circ$ and $-90^\circ$ are the same operation viewed from opposite ends}.

\pq{Methane and $A_4$.}
Two points of $T_+$ differ in exactly two coordinates, so they are $2\sqrt2$ apart, joined by a diagonal of a face. A rotation preserves distances, so $g(T_+)$ is a set of four vertices pairwise $2\sqrt2$ apart. The vertices $2\sqrt2$ from a given vertex are those that differ from it in two coordinates, which have the same product of coordinates. So the four vertices of $g(T_+)$ have equal products: \ans{$g(T_+)$ is $T_+$ (product $+1$) or $T_-$ (product $-1$)}, and $g(T_-)$ is the other. Each $g$ therefore permutes the pair $\{T_+, T_-\}$, exchanging the two exactly when $\chi(g) = -1$, so $\chi(gh) = \chi(g)\chi(h)$.

$R(1, 1, 1) = (-1, 1, 1) \in T_-$, so \ans{$\chi(R) = -1$}, and \ans{$\chi(R^2) = (-1)^2 = +1$}. $C$ permutes the coordinates cyclically and keeps their product: \ans{$\chi(C) = +1$}. $B(1, 1, 1) = (1, 1, -1) \in T_-$: \ans{$\chi(B) = -1$}. These agree with $\sgn(1234) = -1$, $\sgn\bigl((13)(24)\bigr) = +1$, $\sgn(243) = +1$ and $\sgn(13) = -1$. Both $\chi$ and $\sgn \circ \varphi$ are homomorphisms into the Abelian group $\{\pm 1\}$, hence constant on conjugacy classes, and $e, R, R^2, C, B$ represent all five classes by (d). So \ans{$\chi = \sgn \circ \varphi$}. The rotations that map the methane molecule to itself form $\ker\chi = \varphi^{-1}(A_4)$, \ans{a group isomorphic to $A_4$}. The sign of $\varphi(g)$ is \ans{$+1$ when $g$ maps each of the two tetrahedra inscribed in the cube to itself, and $-1$ when it exchanges them}.

\pq{The axes, $\ker\psi$, and $S_4/V \cong S_3$.}
$R$ sends the $x$-axis to the $y$-axis and the $y$-axis to the $x$-axis and fixes the $z$-axis: \ans{$\psi(R) = (xy)$}. $C$ sends $(1, 0, 0) \mapsto (0, 1, 0)$, $(0, 1, 0) \mapsto (0, 0, 1)$, $(0, 0, 1) \mapsto (1, 0, 0)$: \ans{$\psi(C) = (xyz)$}. A transposition and a three-cycle generate $S_3$, so $\psi$ is onto, and the first isomorphism theorem gives \ans{$|\ker\psi| = 24/6 = 4$}. The identity and the three half-turns about the coordinate axes, such as $(x, y, z) \mapsto (x, -y, -z)$, map every axis to itself; that is four elements already, so \ans{$\ker\psi$ consists of $e$ and the three face half-turns}. By (d), $\varphi$ maps them onto $V = \{e, (12)(34), (13)(24), (14)(23)\}$. A kernel is normal, and $\varphi$ is an isomorphism, so \ans{$V \trianglelefteq S_4$ and $S_4/V \cong G/\ker\psi \cong S_3$}. $\ker\psi$ is \ans{$D_2 \cong \Z_2 \times \Z_2$, the group of half-turns of a box} from Week 1.
\end{parts}
%END-SOLUTIONS

\end{document}
